\subsection{CONTINUITY OF INTEGRALS AS FUNCTIONS OF A PARAMETER}

Prop. 6 (IV, p.~171) shows that when a differential equation

\[
\mathbf {x} ^ {\prime} = \mathbf {f} (t, \mathbf {x})
\]

is ``close'' to a Lipschitz equation \(\mathbf{x}^{\prime} = \mathbf{g}(t, \mathbf{x})\) and when one supposes that both equations have an approximate solution on the same interval, then these approximate solutions are ``close''; we shall clarify this result by showing that the existence of solutions of the Lipschitz equation \(\mathbf{x}^{\prime} = \mathbf{g}(t, \mathbf{x})\) on an interval implies that of approximate solutions of \(\mathbf{x}^{\prime} = \mathbf{f}(t, \mathbf{x})\) on the same interval, so long as, on the latter, the values of the solution of \(\mathbf{x}^{\prime} = \mathbf{g}(t, \mathbf{x})\) are not ``too close'' to the boundary of H.

PROPOSITION 7. Let f and g be two functions defined on I × H, satisfying the hypotheses of lemma 1 of IV. p.~165, and such that, on I × H

\[
\left\| \mathbf {f} (t, \mathbf {x}) - \mathbf {g} (t, \mathbf {x}) \right\| \leqslant \alpha .\tag{16}
\]

Suppose further that \(\mathbf{g}\) is Lipschitz with respect to a constant \(k > 0\) on \(\mathrm{I} \times \mathrm{H}\) and that \(\mathbf{f}\) is locally Lipschitz on \(\mathrm{I} \times \mathrm{H}\) , or that \(\mathrm{E}\) is finite dimensional. Let \((t_0, \mathbf{x}_0)\) be a point of \(\mathrm{I} \times \mathrm{H}\) , \(\mu\) a number \(>0\) , and

\[
\varphi (t) = \mu e ^ {k (t - t _ {0})} + \alpha \frac {e ^ {k (t - t _ {0})} - 1}{k}.
\]

Let \(\mathbf{u}\) be an integral of the equation \(\mathbf{x}' = \mathbf{g}(t, \mathbf{x})\) defined on an interval \(\mathrm{K} = [t_0, b[\) contained in I, equal to \(\mathbf{x}_0\) at the point \(t_0\) , and such that for all \(t \in \mathrm{K}\) the closed ball with centre \(\mathbf{u}(t)\) and radius \(\varphi(t)\) is contained in H. Under these conditions, for every \(\mathbf{y} \in \mathrm{H}\) such that \(\| \mathbf{y} - \mathbf{x}_0 \| \leqslant \mu\) there exists an integral \(\mathbf{v}\) of \(\mathbf{x}' = \mathbf{f}(t, \mathbf{x})\) , defined on K, with values in H, and equal to \(\mathbf{y}\) at the point \(t_0\) ; moreover, \(\| \mathbf{u}(t) - \mathbf{v}(t) \| \leqslant \varphi(t)\) on K.

Let M be the family of integrals of \(\mathbf{x}^{\prime} = \mathbf{f}(t, \mathbf{x})\) each of which takes its values in H, is equal to y at \(t_{0}\) , and is defined on a half-open interval \([t_{0}, \tau[\) contained in I (depending on the interval considered). By th. 1 of IV, p.~171 (when f is locally Lipschitz) or IV, p.~167 corollary (when E is finite dimensional), M is not empty, and the same reasoning as in prop. 3 of IV, p.~166 shows that M is inductive for the order ``v is a restriction of w''. Let \(v_{0}\) be a maximal element of M and \([t_{0}, t_{1}[\) the interval of definition of \(v_{0}\) ; by prop. 6 of IV, p.~171, it all comes down to proving that \(t_{1} \geqslant b\) . If not, one would have

\[
\left\| \mathbf {u} (t) - \mathbf {v} _ {0} (t) \right\| \leqslant \varphi (t)
\]

on the interval \([t_0, t_1[\) by prop. 6; now on the compact interval \([t_0, t_1]\) the regulated function \(\mathbf{g}(t, \mathbf{u}(t))\) is bounded, so the function \(\mathbf{g}(t, \mathbf{v}_0(t))\) is bounded on the interval \([t_0, t_1[\) , for \(\| \mathbf{g}(t, \mathbf{v}_0(t)) \| \leqslant \| \mathbf{g}(t, \mathbf{u}(t)) \| + k\varphi(t)\) on this interval. Since \(\mathbf{v}_0\) is an approximate solution of \(\mathbf{x}' = \mathbf{g}(t, \mathbf{x})\) to within \(\alpha\) on \([t_0, t_1[\) there exists a number \(M > 0\) such that \(\| \mathbf{v}_0'(t) \| \leqslant M\) on this interval, except at the points of a countable set; the mean value theorem now shows that \(\| \mathbf{v}_0(s) - \mathbf{v}_0(t) \| \leqslant M |s - t|\) for every pair of points \(s\) , \(t\) of \([t_0, t_1[\) , so (by Cauchy's criterion) \(\mathbf{v}_0(t)\) has a left limit \(\mathbf{c}\) at the point \(t_1\) , and, by continuity, one has \(\| \mathbf{c} - \mathbf{u}(t_1) \| \leqslant \varphi(t_1)\) , and so \(\mathbf{c} \in H\) . Now one sees, from IV, p.~171, th. 1 or IV, p.~167, corollary, that there exists an integral of \(\mathbf{x}' = \mathbf{f}(t, \mathbf{x})\) defined on an interval \([t_1, t_2[\) and equal to \(\mathbf{c}\) at \(t_1\) , which contradicts the definition of \(\mathbf{v}_0\) .

THEOREM 3. Let F be a topological space and let f be a map from I × H × F into E such that, for every \(\xi \in F\) , the function \((t, \mathbf{x}) \mapsto \mathbf{f}(t, \mathbf{x}, \xi)\) is Lipschitz on I × H, and such that, when \(\xi\) tends to \(\xi_{0}\) , \(\mathbf{f}(t, \mathbf{x}, \xi)\) tends uniformly to \(\mathbf{f}(t, \mathbf{x}, \xi_{0})\) on I × H. Let \(\mathbf{u}_{0}(t)\) be an integral of \(\mathbf{x}' = \mathbf{f}(t, \mathbf{x}, \xi_{0})\) defined on an interval J = {[}t₀, b{[} contained in I, with values in H, and equal to \(x_{0}\) at t₀. For every compact interval {[}t₀, t₁{]} contained in J there exists a neighbourhood V of \(\xi_{0}\) in F such that, for every \(\xi \in V\) , the equation \(\mathbf{x}' = \mathbf{f}(t, \mathbf{x}, \xi)\) has an integral (and only one) \(\mathbf{u}(t, \xi)\) defined on {[}t₀, t₁{]}, with values in H and equal to \(x_{0}\) at t₀; moreover, when \(\xi\) tends to \(\xi_{0}\) the solution \(\mathbf{u}(t, \xi)\) tends uniformly to \(\mathbf{u}_{0}(t)\) on {[}t₀, t₁{]}.

Indeed, let r \textgreater{} 0 be such that for \(t_{0} \leqslant t \leqslant t_{1}\) the closed ball with centre \(\mathbf{u}_{0}(t)\) and radius r is contained in H; if \(\mathbf{f}(t, \mathbf{x}, \xi)\) is Lipschitz with respect to the constant k \textgreater{} 0 on I × H we take \(\alpha\) small enough that \(\alpha \frac{e^{k(t_{1}-t_{0})} - 1}{k} < r\) ; taking V such that, for every \(\xi \in V\) , one has \(\|f(t, \mathbf{x}, \xi) - f(t, \mathbf{x}, \xi_{0})\| \leqslant \alpha\) on I × H, one can apply prop. 7 of IV, p.~174; moreover,

\[
\| \mathbf {u} (t, \xi) - \mathbf {u} _ {0} (t) \| \leqslant \alpha \frac {e ^ {k \left(t _ {1} - t _ {0}\right)} - 1}{k}
\]

on \([t_{0}, t_{1}]\) , which completes the proof of the theorem.

Remark. When H = E and the condition (16) of IV, p.~174 is satisfied on I × E, prop. 7 of IV, p.~174 applies to every solution u of \(\mathbf{x}' = \mathbf{g}(t, \mathbf{x})\) on any interval where this solution is defined; one may even assume that \(\mathbf{g}(t,\mathbf{x})\) is Lipschitz with respect to a regulated function \(k(t)\) though not necessarily bounded on this interval.

